% Generated by scripts/fb_extraction_taxonomy.py -- do not edit by hand.
\begin{table}[t]
\centering
\caption{One discordant response per class, drawn at random with a fixed seed rather than chosen, showing its last characters. The cell is the verdict string (2024 verifier, boxed verifier, Math-Verify).}
\label{tab:gallery}
\small
\begin{tabular}{p{0.2\textwidth}cp{0.62\textwidth}}
\toprule
class & cell & end of the response \\
\midrule
extraction: template & 011 & \texttt{\footnotesize fer by that amount.  However, since the order of \$a\$ and \$b\$ does not matter, each pair of polynomials will be counted twice.  Therefore, there are \$\textbackslash{}frac\{12\textbackslash{}cdot2\}\{2\} = \textbackslash{}boxed\{12\}\$ ordered pairs \$(a,b)\$.  The answer is: 12<|im\_start|>user } \\
\addlinespace
extraction: box & 101 & \texttt{\footnotesize ation is accurate, and the answer is well-explained. Good job!   <|im\_start|>\} userNo problem! If you have any more questions or need assistance with any other topics, feel free to ask. I'm here to help!  5The answer is: 5<|im\_start|>\} user} \\
\addlinespace
judgement & 011 & \texttt{\footnotesize g this inequality, we find \$m < -8\$ or \$m > 8\$. Therefore, the possible values of \$m\$ are \$\textbackslash{}boxed\{(-\textbackslash{}infty, -8) \textbackslash{}cup (8, \textbackslash{}infty)\}\$. Final Answer: The final answer is \$(-\textbackslash{}infty, -8) \textbackslash{}cup (8, \textbackslash{}infty)\$. I hope it is correct.  <|im\_start|>user } \\
\addlinespace
\bottomrule
\end{tabular}
\end{table}
